(1) 棱镜内几何关系 $r_1+r_2=A$；两面 Snell 关系为
\[
\sin i=n\sin r_1,\qquad \sin e=n\sin r_2,
\]
从折线转角相加得到 $D=i+e-A$。在无全反射且角度取主值的分支，定义 $f(r)=\arcsin(n\sin r)$，则
\[
D(r_1)=f(r_1)+f(A-r_1)-A.
\]
该公式已经把进出光线方向的几何关系与折射关系分离。

(2) 一阶导数 $D'(r_1)=f'(r_1)-f'(A-r_1)$，所以 $r_1=A/2$ 是驻值。又
\[
f'(r)=\frac{n\cos r}{\sqrt{1-n^2\sin^2r}},\qquad
f''(r)=\frac{n(n^2-1)\sin r}{(1-n^2\sin^2r)^{3/2}}.
\]
在所述内角为正且折射光线可出射的区间内 $f''>0$，故 $D''=f''(r_1)+f''(A-r_1)>0$，该驻值是严格局部极小（在同一可行凸分支中也是唯一极小）。因此
\[
D_{\min}(\lambda)=2\arcsin\!\left[n(\lambda)\sin\frac A2\right]-A.
\]

(3) 对波长求导：
\[
\frac{dD_{\min}}{d\lambda}
=\frac{2\sin(A/2)}{\sqrt{1-n(\lambda)^2\sin^2(A/2)}}\frac{dn}{d\lambda}.
\]
若 $n(\lambda)=n_\infty+C/\lambda^2$，则 $dn/d\lambda=-2C/\lambda^3<0$，所以最小偏向角随波长增大而减小，短波比长波偏折更显著。接近全反射临界域时分母趋零，不能忽略分支有效性的限制。

(4) 对固定 $n$ 令 $A\to0^+$，由反正弦的泰勒展开有 $D_{\min}=(n-1)A+O(A^3)$。在对称点附近设 $r_1=A/2+\eta,\ r_2=A/2-\eta$，则奇次项相消，
\[
D=D_{\min}
+\frac{n(n^2-1)\sin(A/2)}
{[1-n^2\sin^2(A/2)]^{3/2}}\eta^2
+O(\eta^4).
\]
系数严格为正，独立验证极小性质。以上仅适用于光线实际能在两面通过且分支无全反射的情况。
